\begin{papercard}{habib2020}
\cf{Stages} \stagetag{W} (also \stagetag{C} \stagetag{I}). An in-lab usability study ($n=24$) of real websites' opt-out and deletion choices, with tasks drawn from an expert review of 150 sites (Habib et al.\ 2019). Findings are organised by the User Action Framework: finding, learning, using, understanding.

\cf{Aim} A holistic usability evaluation of the \emph{end-to-end interaction} needed to use common implementations of privacy choices, because availability does not imply usability \pc{habib2020}{1}{However, availability does not imply usability, leaving open the question}. Choice types: email marketing (CAN-SPAM), targeted ads, data deletion (GDPR Art.\ 17) \pc{habib2020}{2}{the GDPR grants a “right to be forgot}.

\cf{Task selection} From the earlier 150-site review, three common presentation patterns: account setting, link from the privacy policy, or text instructions in the policy \pc{habib2020}{3}{these choices are most commonly presented using one of three patterns}. Websites using only one pattern were selected, three per task type, giving nine live websites (Table 1 lists the mechanism and the \emph{minimum number of user actions}) \pc{habib2020}{4}{For each of the task types, we selected three websites that fol}.

\cf{Design} Each session had a semi-structured interview on expectations (median 11 minutes), two think-aloud tasks, and follow-up questions \pc{habib2020}{3}{The first portion of the study session, a semi-structured in}. Tasks were counterbalanced and stratified so each participant did two different types, one account-settings and one policy mechanism, giving 12 groupings with pairs in inverse order \pc{habib2020}{4}{To minimize learning effects and prevent fatigue, we counter}. Setup: a fresh Chrome Incognito window, a study Gmail account, and study credentials on index cards, so no personal data was used \pc{habib2020}{4}{Participants were not required to use their own credentials or personal information}. Scenarios were realistic, e.g.\ ``You're uncomfortable with [website] keeping a record of your location'' \pc{habib2020}{4}{and want to remove all of your data from the company’s databases}.

\cf{Data collection} One moderator and one note-taker, with audio and screen recordings under IRB approval \pc{habib2020}{4}{One researcher moderated all participant sessions}. Local recruitment with a screening survey on privacy practices (Pew items) to build a diverse sample, \$20 compensation, median session 50 minutes. A hint was given after eight minutes, a threshold set in pilots, and any assistance was recorded in the analysis \pc{habib2020}{4}{This threshold of eight min}.

\cf{Analysis} Transcripts were auto-transcribed after removing identifying details \pc{habib2020}{4}{We took extra measures to preserve participants}. Interviews were coded inductively. One researcher open-coded to build a codebook; two researchers coded a random six interviews, reviewed disagreements and revised; the remainder were double-coded with disagreements reconciled \pc{habib2020}{5}{Two researchers then collaboratively revised the codebook after individually coding}. Task videos: an analysis template counted interactions and errors, with one researcher coding and another confirming \pc{habib2020}{5}{We created an analysis template to systematically count the interactions and errors}. Findings follow the User Action Framework (planning, translation, physical action, assessment), mapped to \emph{finding, learning, using, understanding} \pc{habib2020}{5}{and mapped these phases of the interaction to finding, learning, using, and understanding}. Frequency words (``a few'', ``some'', ``most'') are defined in Figure 1 instead of exact counts \pc{habib2020}{5}{we adopted the terminology presented in Figure 1}.

\cf{Findings} \emph{Finding}: almost all participants needed help to find account-setting or policy mechanisms (``It's a scavenger hunt'') \pc{habib2020}{6}{Almost all participants required assistance finding the account setting or privacy policy mechanism}. Poor labels misled (``alert preferences'', an ``advertisers'' page), while multiple paths helped \pc{habib2020}{7}{This suggests that more descriptive labels on these websites would help users}. \emph{Learning}: poor formatting misled, e.g.\ data rights that looked clickable, and a Google Analytics opt-out mistaken for an ad opt-out \pc{habib2020}{7}{the opt-out extension only opts users out of Google’s tracking for analytics purposes}. \emph{Using}: on average 26.1 actions for account settings, 37.5 for policy links and 57.6 for policy text, against ideal paths of 9, 22.3 and 9.3 \pc{habib2020}{8}{Using a choice mechanism in account settings resulted in an average of 26.1 user actions}. Written requests lacked needed information, and an unsaved setting caused a post-completion error \pc{habib2020}{9}{This is an example of a post-completion error}. \emph{Understanding}: skepticism that choices are honoured persisted even after deletion \pc{habib2020}{9}{skepticism of deletion choices persisted even after participants used deletion mechanisms}. Ad opt-outs were misunderstood across browsers and devices.

\cf{Limitations stated} Exploratory, with no frequency claims. A small, tech-savvy sample. Nine sites not representative of all categories. Live-site peculiarities. Possible learning across tasks, and priming by the interview \pc{habib2020}{5}{The exploratory nature of this study provides insights into possible usability issues}.

\cf{Recommendations} Unified settings in a standard place, multiple consistent paths and in-context controls, less user effort (toggles instead of written requests), clear labels, and confirmation of when a choice takes effect \pc{habib2020}{10}{Websites should also provide confir}.

\cf{Use for this thesis} Strong precedent for task design and for coding observable walkthrough outcomes. (1) Derive tasks from a prior systematic review of mechanisms, as Habib does from the 150-site review. (2) Use \emph{realistic scenario prompts} and study-only credentials. (3) Count \emph{actions against the ideal path} as an objective effort metric for each privacy task. (4) Code by the finding--learning--using--understanding phases, which align with the CW questions. (5) The coding procedure (open codebook, joint pilot on a random subset, then full double coding with reconciliation) is a citable template, though it is reconciliation, not a reliability statistic. Lab users differ from an expert walkthrough, so transfer the task and metric design, not the user findings.
\end{papercard}
